% =============================================================================
% CHƯƠNG 4 — PHÂN TÍCH KẾT QUẢ NGHIÊN CỨU
% Văn phong khoa học, thuật ngữ và cấu trúc đã được chuẩn hóa.
% =============================================================================

\chapter{PHÂN TÍCH KẾT QUẢ NGHIÊN CỨU}
\label{chap:ket-qua}

\section{Đặc điểm dữ liệu và kết quả hiệu chỉnh}
\label{sec:c4-data-calibration}

\begin{table}[H]
\centering
\small
\caption{Thống kê lợi suất log ngày trong giai đoạn hiệu chỉnh 2005--2019}
\label{tab:c4-descriptive}
\begin{tabular}{@{}lccc@{}}
\toprule
& \textbf{AAPL} & \textbf{JPM} & \textbf{XOM}\\
\midrule
Độ biến động năm & 0,3232 & 0,3729 & 0,2341\\
Kurtosis & 9,02 & 21,78 & 16,50\\
Skewness & $-0,30$ & $-0,27$ & $-0,13$\\
\midrule
\multicolumn{4}{l}{\textit{Tương quan cặp}}\\
AAPL--JPM & \multicolumn{3}{c}{0,399}\\
AAPL--XOM & \multicolumn{3}{c}{0,377}\\
JPM--XOM & \multicolumn{3}{c}{0,468}\\
\bottomrule
\end{tabular}
\end{table}

Kurtosis của AAPL, JPM và XOM lần lượt bằng 9,02, 21,78 và 16,50, cao hơn mức
3 của phân phối Gaussian. Skewness của cả ba chuỗi mang dấu âm. Tương quan cặp
nằm trong khoảng 0,377--0,468. Các giá trị tương quan được tính trên toàn bộ
giai đoạn hiệu chỉnh và không biểu diễn biến động của tương quan theo thời
gian.

SBTS sử dụng $(h^{*},K^{*})=(0{,}030,5)$. GBM sử dụng độ biến động và ma trận
tương quan mẫu. Heston sử dụng phương sai dài hạn riêng cho từng tài sản và
các tham số động lực chung theo đặc tả tại Chương~\ref{chap:phuong-phap}. Kết
quả được diễn giải theo đúng ba cấu hình hiệu chỉnh này.

\section{Đối chiếu phân phối của ba mô hình sinh dữ liệu}
\label{sec:c4-distribution}

\begin{table}[H]
\centering
\small
\caption{Độ trung thực phân phối của các quỹ đạo tổng hợp so với dữ liệu lịch sử}
\label{tab:c4-distributional-fidelity}
\begin{tabular}{@{}lcccc@{}}
\toprule
\textbf{Chỉ tiêu} & \textbf{Lịch sử} & \textbf{GBM} &
\textbf{Heston} & \textbf{SBTS}\\
\midrule
Kurtosis bình quân & 15,75 & 3,00 & 4,54 & \textbf{15,28}\\
$|\Delta\mathrm{Kurt}|$ & -- & 12,75 & 11,20 & \textbf{0,47}\\
MAE tương quan & -- & \textbf{0,0002} & 0,040 & 0,034\\
KS bình quân & -- & 0,091 & 0,059 & \textbf{0,031}\\
$\widehat\lambda_L$ tại 5\% & 0,311 & 0,197 & 0,175 & \textbf{0,297}\\
$|\Delta\lambda_L|$ & -- & 0,114 & 0,136 & \textbf{0,014}\\
\bottomrule
\end{tabular}
\end{table}

SBTS có sai lệch kurtosis bằng 0,47, KS bình quân bằng 0,031 và sai lệch phụ
thuộc đuôi trái bằng 0,014. Các giá trị tương ứng của GBM là 12,75, 0,091 và
0,114; của Heston là 11,20, 0,059 và 0,136.

GBM có MAE tương quan bằng 0,0002 vì ma trận tương quan mẫu được sử dụng trực
tiếp trong phân rã Cholesky. Heston có kurtosis và KS gần dữ liệu lịch sử hơn
GBM, nhưng hệ số phụ thuộc đuôi trái bằng 0,175, thấp hơn mức 0,197 của GBM và
0,311 của dữ liệu lịch sử.

\section{Hiệu quả hedging tại ba giai đoạn thị trường}
\label{sec:c4-performance}

\subsection{Giai đoạn đại diện (2023--2025)}

\begin{table}[H]
\centering
\scriptsize
\setlength{\tabcolsep}{3pt}
\caption{Độ lệch chuẩn sai số hedging $\sigma(R)$ trong giai đoạn đại diện (2023--2025)}
\label{tab:c4-recent-main}
\begin{tabular}{@{}llccccc@{}}
\toprule
\textbf{Hợp đồng} & $\boldsymbol\kappa$ & \textbf{GBM} & \textbf{Heston} &
\textbf{SBTS} & \textbf{Giảm so với GBM} & $\boldsymbol p_{\mathrm{BH}}$\\
\midrule
Basket call & 0,95 & $0,0059\pm0,0005$ & $0,0056\pm0,0005$ &
$\mathbf{0,0034\pm0,0005}$ & 42\% & $<0,001$\\
Basket call & 1,00 & $0,0076\pm0,0005$ & $0,0074\pm0,0005$ &
$\mathbf{0,0060\pm0,0011}$ & 21\% & 0,004\\
Basket call & 1,05 & $0,0086\pm0,0004$ & $0,0082\pm0,0004$ &
$\mathbf{0,0066\pm0,0008}$ & 22\% & $<0,001$\\
\midrule
Worst-of put & 0,95 & $0,0134\pm0,0010$ & $0,0133\pm0,0012$ &
$\mathbf{0,0116\pm0,0013}$ & 13\% & 0,012\\
Worst-of put & 1,00 & $0,0145\pm0,0010$ & $0,0133\pm0,0014$ &
$\mathbf{0,0112\pm0,0019}$ & 23\% & 0,002\\
Worst-of put & 1,05 & $0,0162\pm0,0012$ & $0,0145\pm0,0010$ &
$\mathbf{0,0133\pm0,0020}$ & 18\% & 0,004\\
\bottomrule
\end{tabular}
\end{table}

Trong giai đoạn đại diện, SBTS có $\sigma(R)$ thấp nhất ở sáu tổ hợp hợp đồng--mức giá thực
hiện. Mức giảm so với GBM nằm trong khoảng 13--42\%; mức giảm so với Heston
nằm trong khoảng 8--38\%. Sáu so sánh SBTS--GBM có
$p_{\mathrm{BH}}<0{,}05$. Năm trong sáu so sánh SBTS--Heston có khác biệt
thống kê; worst-of put tại $\kappa=1{,}05$ chưa có khác biệt thống kê.

Đối với basket call, mức giảm lớn nhất so với GBM bằng 42\% tại
$\kappa=0{,}95$. Đối với worst-of put, mức giảm lớn nhất bằng 23\% tại
$\kappa=1{,}00$. Basket call phụ thuộc vào trung bình của rổ; worst-of put phụ
thuộc vào tài sản có giá trung bình thấp nhất.

\subsection{Giai đoạn phục hồi (2021--2022)}

Trong giai đoạn phục hồi, SBTS có $\sigma(R)$ thấp hơn GBM và Heston tại ba mức giá
thực hiện của basket call, với $p_{\mathrm{BH}}<0{,}05$. Đối với worst-of put,
GBM có $\sigma(R)$ thấp hơn SBTS tại $\kappa=0{,}95$ và $1{,}00$; Heston có
$\sigma(R)$ thấp hơn SBTS tại $\kappa=1{,}05$. Các so sánh còn lại chưa có
khác biệt thống kê sau điều chỉnh BH.

\subsection{Giai đoạn căng thẳng (2019--2020)}

Thứ hạng trong giai đoạn căng thẳng đảo chiều so với giai đoạn đại diện. Ở năm trong sáu tổ hợp của mỗi cặp
SBTS--GBM và SBTS--Heston, mô hình tham số có $\sigma(R)$ thấp hơn với
$p_{\mathrm{BH}}<0{,}05$. Tổ hợp worst-of put tại $\kappa=1{,}05$ có chênh lệch theo
cùng hướng bất lợi cho SBTS nhưng chưa đạt ngưỡng sau điều chỉnh
($p_{\mathrm{BH}}=0{,}09$).

$\sigma(R)$ của cả ba mô hình sinh dữ liệu trong giai đoạn căng thẳng đều cao hơn trong giai đoạn đại diện.
Kết quả chỉ xác định thứ hạng tương đối giữa ba cấu hình. Phân tích không kiểm
định riêng mức độ co của hỗ trợ kernel, nên giới hạn ngoại suy của SBTS được
ghi nhận như một cơ chế lý thuyết và không được xem là quan hệ nhân quả đã xác
định.

\subsection{So sánh tổng hợp tại ba giai đoạn}

\begin{table}[H]
\centering
\small
\caption{Kết quả kiểm định từng cặp trên 18 tổ hợp giai đoạn--hợp đồng--mức giá thực hiện}
\label{tab:c4-scoreboard}
\begin{tabular}{@{}lccc@{}}
\toprule
\textbf{Cặp so sánh} & \textbf{Mô hình thứ nhất thắng} &
\textbf{Mô hình thứ hai thắng} & \textbf{Chưa phân biệt}\\
\midrule
SBTS so với GBM & \textbf{9} & 7 & 2\\
SBTS so với Heston & \textbf{8} & 6 & 4\\
Heston so với GBM & \textbf{3} & 2 & 13\\
\bottomrule
\end{tabular}
\end{table}

SBTS có 9 tổ hợp tốt hơn GBM, 7 tổ hợp kém hơn và 2 tổ hợp chưa có khác biệt
thống kê. So với Heston, SBTS có 8 tổ hợp tốt hơn, 6 tổ hợp kém hơn và 4 tổ
hợp chưa có khác biệt thống kê. Các tổ hợp bất lợi cho SBTS tập trung trong
giai đoạn căng thẳng và nhóm worst-of put của giai đoạn phục hồi.

Cặp Heston--GBM có 13/18 tổ hợp chưa có khác biệt thống kê. Trong đặc tả được
sử dụng, Heston không tạo mức giảm $\sigma(R)$ đồng đều so với GBM.

\section{Kết quả của quy trình MSE--CVaR}
\label{sec:c4-curriculum}

\begin{table}[H]
\centering
\small
\caption{Thay đổi sau giai đoạn CVaR tại giai đoạn đại diện (2023--2025), $\kappa=1{,}00$}
\label{tab:c4-curriculum}
\begin{tabular}{@{}llccr@{}}
\toprule
\textbf{Hợp đồng} & \textbf{Mô hình sinh dữ liệu} & \textbf{Chỉ MSE} &
\textbf{Sau CVaR} & \textbf{Mức giảm}\\
\midrule
Basket call & GBM & 0,0079 & 0,0076 & $+4,7\%$\\
Basket call & Heston & 0,0072 & 0,0074 & $-2,7\%$\\
Basket call & SBTS & 0,0074 & \textbf{0,0060} & \textbf{$+19,3\%$}\\
\midrule
Worst-of put & GBM & 0,0148 & 0,0145 & $+1,9\%$\\
Worst-of put & Heston & 0,0131 & 0,0133 & $-1,3\%$\\
Worst-of put & SBTS & 0,0099 & 0,0112 & $-13,4\%$\\
\bottomrule
\end{tabular}
\end{table}

Dấu dương trong Bảng~\ref{tab:c4-curriculum} biểu thị $\sigma(R)$ giảm sau
giai đoạn CVaR. Đối với basket call, $\sigma(R)$ của SBTS giảm 19,3\%; mức thay
đổi của GBM và Heston lần lượt bằng $+4,7\%$ và $-2,7\%$. Đối với worst-of
put, $\sigma(R)$ của SBTS tăng 13,4\%; mức thay đổi của GBM và Heston lần lượt
bằng $+1,9\%$ và $-1,3\%$. Kết quả của giai đoạn CVaR phụ thuộc vào hợp đồng
và mô hình sinh dữ liệu.

\section{Mức vốn ban đầu $V_0^{*}$}
\label{sec:c4-v0}

\begin{table}[H]
\centering
\small
\setlength{\tabcolsep}{4pt}
\caption{Mức vốn ban đầu $V_0^{*}$, trung bình qua mười seed}
\label{tab:c4-v0}
\begin{tabular}{@{}lcccccc@{}}
\toprule
& \multicolumn{3}{c}{\textbf{Basket call}}
& \multicolumn{3}{c}{\textbf{Worst-of put}}\\
\cmidrule(lr){2-4}\cmidrule(lr){5-7}
\textbf{Mô hình sinh dữ liệu} & 0,95 & 1,00 & 1,05 & 0,95 & 1,00 & 1,05\\
\midrule
GBM & 0,0857 & 0,0587 & 0,0385 & 0,0990 & 0,1372 & 0,1796\\
Heston & 0,0879 & 0,0591 & 0,0372 & 0,1105 & 0,1469 & 0,1881\\
SBTS & \textbf{0,0696} & \textbf{0,0406} & \textbf{0,0210}
& \textbf{0,0738} & \textbf{0,1118} & \textbf{0,1555}\\
\bottomrule
\end{tabular}
\end{table}

SBTS có $V_0^{*}$ thấp nhất tại sáu tổ hợp. Mức giảm so với GBM nằm trong
khoảng 13--45\%; mức giảm so với Heston nằm trong khoảng 17--44\%. Dưới điều
kiện bậc nhất của MSE, $V_0^{*}$ làm sai số hedging kỳ vọng bằng không dưới
phân phối huấn luyện. Đại lượng này còn được gọi là giá bàng quan bậc hai
(\textit{quadratic indifference price}).

$V_0^{*}$ thấp hơn biểu thị mức vốn ban đầu thấp hơn trong cùng điều kiện hợp
đồng và chi phí giao dịch. Kết quả này được xem xét cùng $\sigma(R)$. Trong
giai đoạn đại diện, SBTS có cả $V_0^{*}$ và $\sigma(R)$ thấp hơn hai mô hình tham số. Trong
giai đoạn căng thẳng, thứ hạng $\sigma(R)$ đảo chiều. Vì vậy, $V_0^{*}$ không được diễn giải
tách khỏi kết quả hedging và không được xem là giá thị trường độc lập với
mô hình.

\section{Kết quả Model Confidence Set}
\label{sec:c4-mcs}

\begin{table}[H]
\centering
\scriptsize
\setlength{\tabcolsep}{4pt}
\caption{Thành viên Model Confidence Set ở mức ý nghĩa 10\%}
\label{tab:c4-mcs}
\begin{tabular}{@{}llccc@{}}
\toprule
\textbf{Giai đoạn} & \textbf{Hợp đồng} & $\boldsymbol\kappa=0,95$ &
$\boldsymbol\kappa=1,00$ & $\boldsymbol\kappa=1,05$\\
\midrule
Căng thẳng & Basket call & \{GBM,Heston,SBTS\} & \{GBM,Heston,SBTS\} &
\{GBM,Heston,SBTS\}\\
Căng thẳng & Worst-of put & \{GBM,Heston,SBTS\} & \textbf{\{Heston\}} &
\{GBM,Heston,SBTS\}\\
\midrule
Phục hồi & Basket call & \textbf{\{SBTS\}} & \{GBM,Heston,SBTS\} &
\textbf{\{SBTS\}}\\
Phục hồi & Worst-of put & \{GBM,Heston,SBTS\} & \textbf{\{GBM\}} &
\{GBM,Heston\}\\
\midrule
Đại diện & Basket call & \textbf{\{SBTS\}} & \textbf{\{SBTS\}} &
\textbf{\{SBTS\}}\\
Đại diện & Worst-of put & \textbf{\{SBTS\}} & \textbf{\{SBTS\}} &
\textbf{\{SBTS\}}\\
\bottomrule
\end{tabular}
\end{table}

Trong giai đoạn đại diện, SBTS là thành viên duy nhất của MCS tại sáu tổ hợp. Trong
giai đoạn phục hồi, SBTS là phần tử duy nhất tại hai tổ hợp basket call; GBM là phần tử
duy nhất tại worst-of put với $\kappa=1{,}00$; các tổ hợp còn lại giữ từ hai
đến ba mô hình. Trong giai đoạn căng thẳng, năm tổ hợp giữ cả ba mô hình; worst-of put với
$\kappa=1{,}00$ chỉ giữ Heston.

Tổng cộng, SBTS xuất hiện trong MCS tại 15/18 tổ hợp và là phần tử duy nhất tại
8 tổ hợp. GBM và Heston cùng xuất hiện tại 9/18 tổ hợp; mỗi mô hình là phần tử
duy nhất một lần.

Kiểm định ghép cặp sử dụng chênh lệch trong từng seed. MCS đánh giá đồng thời
ba mô hình và sử dụng phân tán giữa các seed. Vì vậy, một tổ hợp có thể có
khác biệt trong kiểm định ghép cặp nhưng vẫn giữ nhiều mô hình trong MCS.

\section{Thảo luận kết quả}
\label{sec:c4-discussion}

\subsection{Kết quả tại giai đoạn đại diện}

Trên dữ liệu giai đoạn hiệu chỉnh, SBTS có sai lệch kurtosis, KS và sai lệch
phụ thuộc đuôi trái thấp nhất. Trong giai đoạn đại diện, chiến lược được huấn luyện bằng dữ liệu SBTS có $\sigma(R)$ thấp
nhất tại sáu tổ hợp. MCS giữ SBTS là phần tử duy nhất tại cả sáu tổ hợp. Kiến
trúc mạng, số quỹ đạo, hợp đồng, chi phí giao dịch và quy trình tối ưu được giữ
cố định giữa ba mô hình sinh dữ liệu.

\subsection{Sự đảo chiều dưới dịch chuyển phân phối}

Trong giai đoạn phục hồi, thứ hạng của worst-of put khác thứ hạng của basket call.
Trong giai đoạn căng thẳng, thứ hạng của SBTS đảo chiều so với giai đoạn đại diện tại năm trong sáu tổ hợp
của mỗi cặp so sánh. SBTS ước lượng chuyển tiếp từ lân cận thực nghiệm; GBM và
Heston sinh nhiễu từ phân phối tham số. Nghiên cứu chưa đo trực tiếp số lân cận
kernel trong từng cửa sổ đánh giá, nên không phân rã mức đóng góp của miền hỗ
trợ vào sự đảo chiều.

\subsection{Lựa chọn mô hình sinh dữ liệu theo từng giai đoạn}

Trong giai đoạn đại diện, kiểm định từng cặp và MCS đều lựa chọn SBTS. Trong giai đoạn phục hồi,
SBTS được lựa chọn cho basket call, trong khi một số tổ hợp worst-of put lựa
chọn GBM hoặc Heston. Trong giai đoạn căng thẳng, MCS giữ cả ba mô hình tại năm trong sáu tổ
hợp. Vì vậy, thứ hạng được báo cáo riêng theo giai đoạn, hợp đồng và mức giá
thực hiện.

Checkpoint MSE và CVaR cũng tạo kết quả khác nhau. Sau giai đoạn CVaR,
$\sigma(R)$ của basket call SBTS giảm 19,3\%, trong khi $\sigma(R)$ của
worst-of put SBTS tăng 13,4\% tại $\kappa=1{,}00$. Việc lựa chọn mô hình sinh
dữ liệu được thực hiện cùng với việc lựa chọn hàm mục tiêu huấn luyện.

\section{Tóm tắt chương}

Trên dữ liệu giai đoạn hiệu chỉnh, SBTS có sai lệch kurtosis, KS và phụ thuộc
đuôi trái thấp nhất. Trong giai đoạn đại diện,
SBTS có $\sigma(R)$ thấp nhất và là phần tử duy nhất của MCS tại sáu tổ hợp.
Trong giai đoạn phục hồi, kết quả khác nhau giữa basket call và worst-of put. Trong
giai đoạn căng thẳng, mô hình tham số có $\sigma(R)$ thấp hơn SBTS tại năm trong sáu tổ hợp
của mỗi cặp so sánh. SBTS có $V_0^{*}$ thấp nhất tại sáu tổ hợp, nhưng thứ hạng
này được diễn giải cùng kết quả hedging theo từng giai đoạn thị trường.
